% ============================================================
%  4. METHODOLOGY
% ============================================================
\chapter{Methodology}\label{sec:methodology}

% \section{Forecasting Framework}
% All models are estimated on a training sample running from December 2000 to
% March 2024 and evaluated on a hold-out test set covering the final 24 months
% from April 2024 to March 2026. Model selection and specification are conducted
% exclusively on the training sample to avoid look-ahead bias. As a robustness
% check, Time Series Cross-Validation (TSCV) is applied to the ARIMA benchmark
% using an expanding window with a minimum training length of 60 months,
% averaging one-step-ahead forecast errors across the full sample.

\section{ARIMA Models}

\paragraph{Naive ARIMA (\textit{Model 1})}
The first model is a naive univariate ARIMA, forecasting Danish inflation using only its own past values and imposing no economic structure. The general seasonal ARIMA$(p,d,q)(P,D,Q)_{12}$ model is written as:
\begin{equation}
    \Phi_P(B^{12})\,\phi_p(B)\,\Delta^d\Delta^D_{12}\,y_t
= \Theta_Q(B^{12})\,\theta_q(B)\,\varepsilon_t,
\end{equation}
where $B$ is the backshift operator, $\phi_p$ and $\theta_q$ are the non-seasonal AR and MA polynomials, $\Phi_P$ and $\Theta_Q$ are their seasonal counterparts, and $\varepsilon_t$ is white noise. Order selection follows \textcite{box_jenkins_1976} and is automated by minimising AICc over a grid of candidate specifications \parencite{hyndman_athanasopoulos_2021}. 
% Residual adequacy is assessed using the Ljung-Box test.

\paragraph{Pre-Surge ARIMA (\textit{Model 2})}
The second model tests whether pre-surge inflation dynamics retain out-of-sample predictive value. The ARIMA specification is estimated exclusively on data prior to April 2021, fixing all coefficients at their pre-surge values. This directly tests whether the inflation surge represents a structural break or a temporary deviation from an otherwise stable process.

\section{Dynamic Regression Models} \label{sec:dynregmodel}

The dynamic regression models extend the ARIMA benchmark by incorporating
external predictors while preserving the ability to capture serial dependence.
The general specification is:
% \[
% y_t = \symbf{x}_t'\symbfit{\beta} + \eta_t, \qquad
% \phi_p(B)\,\eta_t = \theta_q(B)\,\varepsilon_t,
% \]
\begin{equation}
    y_t = \mathbf{x}'_t\boldsymbol{\beta} + \eta_t, \qquad
\Phi_P(B^{12})\,\phi_p(B)\,\eta_t = \Theta_Q(B^{12})\,\theta_q(B)\,\varepsilon_t,
\end{equation}

where $y_t$ is Danish inflation, $\mathbf{x}_t$ contains the external
regressors, and $\eta_t$ follows an ARIMA$(p,0,q)$ process with order selected
by AIC minimisation.

\paragraph{Univariate Exogenous Regressors (\textit{Model 3})} Tests the domestic Phillips Curve - the hypothesis that lower unemployment drives higher inflation - by regressing Danish inflation on first-differenced Danish unemployment at lags 0 to 6:
\begin{equation}
    y_t = \sum_{j=0}^{6} \gamma_j\,\Delta u^{\text{DK}}_{t-j} + \eta_t.
\end{equation}

Lags 0 to 6 are included to allow for delayed transmission of labour market
conditions to prices, capturing effects up to two quarters ahead.

\paragraph{Multivariate Exogenous Regressors (\textit{Model 4})} Addresses the risk of overfitting inherent in the
pre-specified lag structure of Model 3 by applying General-to-Specific (GETS)
model selection \parencite{doornik2014}. A General Unrestricted Model (GUM)
is specified with six autoregressive lags of Danish inflation and six lags each
of first-differenced Danish unemployment and EU27 inflation as
candidate regressors:
\begin{equation}
    y_t = \mu + \sum_{k=1}^{6} \phi_k\,y_{t-k}
    + \delta\,s_t
    + \sum_{j=0}^{6} \gamma_j\,\Delta u^{\text{DK}}_{t-j}
    + \sum_{j=0}^{6} \lambda_j\, \pi^{\text{EU27}}_{t-j}
    + \varepsilon_t,
\end{equation}
where $s_t$ is a surge dummy equal to one from April 2021 to July 2023 and zero otherwise, included as a selectable regressor subject to elimination. The GUM is estimated by OLS with heteroskedasticity-robust (White) standard errors to account for potential variance instability during the inflation surge. The \texttt{getsm()} algorithm then sequentially eliminates regressors whose $t$-statistic falls below the ten percent critical value, checking at each step that residual white noise is preserved. The retained specification defines Model~4. The inclusion of EU27 inflation as a candidate regressor is motivated by the Granger causality results in Section~\ref{sec:granger}.

\section{Forecast Evaluation}
Forecast accuracy is evaluated using the Root Mean Squared Error, $\text{RMSE} = \sqrt{\frac{1}{h}\sum_{t=1}^{h}(y_t - \hat{y}_t)^2}$, and
Mean Absolute Error, $\text{MAE}  = \frac{1}{h}\sum_{t=1}^{h}|y_t - \hat{y}_t|$. Unless stated otherwise, $h = 24$ is the forecast horizon and $\hat{y}_t$ is the model forecast. RMSE penalises large errors more heavily due to the quadratic term, while MAE treats all errors equally. A model is considered to improve upon the ARIMA benchmark if it achieves a lower RMSE and MAE over the test period. For the multivariate models (3 and 4), forecasts are conditional on actual observed values of the external regressors in the test period, isolating the incremental value of each predictor channel from any regressors forecasting error.

\subsection{Ljung--Box Test}
Residual adequacy is assessed via the Ljung-Box test. Under the null hypothesis of no autocorrelation up to lag $\ell$, the test 
statistic

\vspace{-0.5cm}
\begin{equation}
    Q^{*} = T(T+2)\sum_{k=1}^{\ell}\frac{\hat{\rho}_k^2}{T-k}
\end{equation}

is asymptotically $\chi^2(\ell)$ distributed, where $T$ is the number of observations and $\hat{\rho}_k$ is the sample autocorrelation at lag $k$. Failure to reject indicates that residuals are consistent with white noise, confirming well-specified dynamics.